% ---
% §1.2 Planteamiento del problema. Reemplaza las líneas 573–603 de Docs/Technical_Report.tex.
% Pregunta de investigación reformulada con el autor el 2026-10-15 (D-041): neutral («modifican»,
% «hasta qué punto»), «características radiómicas» en lugar de «morfológicas» y comparación
% anclada en el mismo vector de características, no en la dimensión.
% Corrige de la auditoría (Docs/insumos/auditoria_s1_s2.md): paráfrasis sin cita de
% Álvarez-Estévez (l. 587 y 591), «geometría euclidiana insuficiente» (refutado por los
% resultados), «caídas de AUC de hasta 0.44 puntos» (Wang et al. reportan AUC de 0.44–0.65,
% no una caída de 0.44) y la presuposición de mejora en la pregunta.
% Decisiones y hallazgos que lo respaldan: D-019, D-031, D-032; H-031, H-039, H-042.
% ---

\subsection{Planteamiento del problema}
\label{sec:planteamiento}

\textbf{La representación como problema}

En el prediagnóstico asistido por computadora, el clasificador no trabaja con la mamografía, sino
con una representación de cada lesión: un vector de descriptores cuantitativos de intensidad,
textura y forma extraídos de la región de interés. La calidad de esa representación limita lo que
cualquier clasificador puede aprender de ella. En el CBIS-DDSM, además, las clases no se distinguen
con la misma facilidad en los dos tipos de hallazgo: en las masas, los descriptores del radiólogo
separan bien las formas irregulares y espiculadas de las ovaladas y circunscritas, mientras que en
las calcificaciones el tipo más frecuente, el pleomórfico, aparece en ambas clases
(capítulo~\ref{sec:eda}). Cabe preguntarse, entonces, si una transformación de esos descriptores
puede reorganizarlos de modo que las lesiones benignas y malignas queden mejor separadas, y si esa
reorganización ayuda a clasificarlas.

\textbf{La propuesta de las codificaciones cuánticas}

Un circuito de codificación $U(\mathbf{x})$ asocia a cada vector de características $\mathbf{x}$ un
estado $|\phi(\mathbf{x})\rangle$ en un espacio de Hilbert de dimensión $2^n$ para $n$ qubits, y la
similitud entre dos lesiones pasa a medirse por la fidelidad
$|\langle\phi(\mathbf{x})|\phi(\mathbf{x}')\rangle|^2$ entre sus estados
(sección~\ref{sec:feature_map}). Los mapas que incluyen compuertas de entrelazamiento, como el
mapa ZZ introducido por Havlíček et al. \cite{havlicek2019supervised}, generan términos que dependen
de pares de características y se han propuesto bajo la conjetura de una posible ventaja cuántica a
corto plazo \cite{AlvarezEstevez2024QuantumKernelBench}. Esa es la motivación de emplearlos como
transformación de los descriptores: no buscar un espacio más grande por sí mismo, sino una geometría
de similitud distinta de la que inducen las transformaciones clásicas.

% respaldo: «If the geometric difference is small, then a classical ML method is guaranteed to provide similar or better performance in prediction on the dataset, independent of the function values or labels.» (huang2021power, p. 2)
% respaldo: «The small geometric difference is a consequence of the exponentially large Hilbert space employed by existing quantum models, where all inputs are too far apart.» (huang2021power, p. 2)
% respaldo: «We show that, under certain conditions, values of quantum kernels over different input data can be exponentially concentrated (in the number of qubits) towards some fixed value.» (Thanasilp2024, p. 1)
% respaldo: «Thus on training with a polynomial number of measurements, one ends up with a trivial model where the predictions on unseen inputs are independent of the input data.» (Thanasilp2024, p. 1)
% respaldo (resumen de arXiv, sin PDF; Docs/fuentes/meta/AlvarezEstevez2024.xml): "However, the practical applicability of these methods remains an open question, particularly beyond the context of specifically-crafted toy problems, and given the current limitations of quantum hardware."
La teoría, sin embargo, no garantiza que esa geometría sea más útil. Huang et al.
\cite{huang2021power} muestran que, si la diferencia geométrica entre un kernel cuántico y los
clásicos es pequeña, un método clásico ofrece un desempeño similar o mejor con cualquier
etiquetado, y observan que en los modelos cuánticos habituales esa diferencia es pequeña porque,
en un espacio de Hilbert exponencialmente grande, todas las entradas quedan demasiado lejos entre
sí. Thanasilp et al. \cite{Thanasilp2024} demuestran que, bajo ciertas condiciones, los valores de
un kernel cuántico se concentran exponencialmente en el número de qubits, de modo que estimarlos
con un número polinomial de mediciones produce un modelo cuyas predicciones no dependen de los
datos. Que una codificación cuántica separe mejor las clases sobre datos clínicos reales es, por
tanto, una pregunta empírica: la aplicabilidad práctica de estos métodos más allá de problemas de
prueba construidos a propósito, y con las limitaciones actuales del \textit{hardware}, sigue
abierta \cite{AlvarezEstevez2024QuantumKernelBench}.

\textbf{La brecha en mamografía}

% respaldo: «In addition, in this research we conduct the experiments using the angle embedding technique since with this method the features can be learned.» (Azevedo2022QuantumTransfer, p. 12)
% respaldo: «As future work, an interesting approach could be to separate the dataset into masses and calcifications and only then perform classification.» (Azevedo2022QuantumTransfer, p. 12)
% respaldo: «Lastly, the best way to know how well this model performs would be to see how the data appears in our Hilbert space. Ideally, both classes would appear very separately in it, as tight clusters.» (Azevedo2022QuantumTransfer, p. 12)
En mamografía, el antecedente más cercano con circuitos cuánticos es el trabajo de Azevedo et al.
\cite{Azevedo2022QuantumTransfer}, que emplean una red híbrida con codificación angular sobre el
BCDR. Entre las líneas que sus autores dejan como trabajo futuro hay dos que este trabajo retoma como ejes: separar masas y
calcificaciones antes de clasificar, y observar cómo se distribuyen los datos en el espacio de
Hilbert, idealmente como conglomerados compactos y separados. La revisión del
capítulo~\ref{sec:antecedentes} no identificó trabajos que evalúen ninguna de las dos sobre el
CBIS-DDSM.

\textbf{La elección del mapa de características}

% respaldo (resumen de arXiv, sin PDF): "Despite outperforming classical methods in ad-hoc datasets, mixed results are obtained for the general case among standard classical benchmarks." / "Instead, it is suggested that a careful choice of the quantum feature map in connection with proper hyperparameterization may prove more effective."
% fuente: D-031
El desempeño de los métodos de kernel cuánticos depende del mapa de características. En una
comparación sobre conjuntos de datos de referencia, Álvarez-Estévez obtiene resultados mixtos
frente a los métodos clásicos y sugiere que una elección cuidadosa del mapa, con una
hiperparametrización adecuada, puede ser más efectiva que optimizar el kernel
\cite{AlvarezEstevez2024QuantumKernelBench}. Este trabajo compara dos mapas que comparten el
acoplamiento ZZ entre pares de características y difieren solo en el eje de la codificación de
primer orden: \textit{ZZFeatureMap}, que gira cada qubit alrededor del eje Z, y
\textit{PauliFeatureMap} con los términos Y y ZZ, que lo gira alrededor del eje Y
(sección~\ref{sec:modulo5}).

\textbf{Formulación del problema de investigación}

A partir de lo anterior, el problema de investigación del presente trabajo se formula de la
siguiente manera:

\begin{quote}
  ¿En qué medida las transformaciones cuánticas variacionales implementadas mediante
  \textit{ZZFeatureMap} y \textit{PauliFeatureMap} en Qiskit modifican la separabilidad de clases
  benigno/maligno de características radiómicas extraídas de mamografías del CBIS-DDSM, y hasta qué
  punto esa diferencia se refleja en la capacidad de clasificación de una misma red neuronal
  clásica, en comparación con representaciones clásicas construidas a partir del mismo vector de
  características?
\end{quote}

La pregunta tiene tres dimensiones, que se abordan de forma integrada. La primera es la de la
representación: si la separabilidad de clases cambia al pasar de las representaciones clásicas a
las cuánticas, medida con métricas geométricas y de kernel (objetivo específico 3). La segunda es
la de la clasificación: si las diferencias de separabilidad se reflejan en el desempeño de un mismo
clasificador y en el orden en que quedan las condiciones (objetivo específico 5). La tercera
acompaña a las dos anteriores y es la de la viabilidad: qué restricciones impone simular estos
circuitos en una computadora clásica y qué implicarían en un procesador cuántico (objetivo
específico 6). La pregunta no presupone el sentido de la respuesta: una separabilidad igual o menor
que la de las representaciones clásicas es también un resultado válido. Los criterios con que se
responde cada dimensión se fijan en el capítulo~\ref{sec:metodologia}.
